% !TEX program = xelatex
\documentclass{英语试卷}
\begin{document}
\examtitle{2024 年 11 月稽阳联谊学校高三联考}{英语}
\par 英语学科 试题卷\par
\par 命题人：嵊州中学 陈淑君 新昌中学 吴燕燕 萧山中学 周旋 审稿人：诸暨中学 钟玲萍 注意事项：\par
\par 1. 答卷前，考生务必将自己的姓名、准考证号填写在答题卡上。\par
\par 2. 回答选择题时，选出每小题答案后，用铅笔把答题卡上对应题目的答案标号涂黑。如需改动，用橡皮擦干净后，再选涂其他答案标号。回答非选择题时，将答案写在答题卡上，写在本试卷上无效。\par
\par 3. 考试结束后，将本试卷和答题卡一并交回。\par
\parttitle{第一部分 听力（共两节，满分 30 分）}
\par 做题时，先将答案标在试卷上，录音内容结束后，你将有两分钟的时间将试卷上的答案转涂到答题卡上。\par
\sectitle{第一节（共5小题：每小题1.5分，满分7.5分）}
\par 听下面 5 段对话。每段对话后有一个小题，从题中所给的 A、B、C 三个选项中选出最佳选项。听完每段对话后，你都有 10 秒钟的时间来回答有关小题和阅读下一小题。每段对话仅读一遍。\par
\question{1. Where does the conversation probably take place?}
\begin{choices}(3)
  \item On the street.
  \item In a café.
  \item In a bakery.
\end{choices}
\question{2. What are the speakers mainly talking about?}
\begin{choices}(3)
  \item A new program.
  \item A speech.
  \item An appointment.
\end{choices}
\question{3. When is the woman's reservation?}
\begin{choices}(3)
  \item On October 16th .
  \item On October 18th .
  \item On October 25th .
\end{choices}
\question{4. What does the man suggest the woman do?}
\begin{choices}(3)
  \item Use an app.
  \item Join a group.
  \item Bring a book.
\end{choices}
\question{5. Whose advice does the woman want to seek?}
\begin{choices}(3)
  \item Mr. Ward's.
  \item Mr. Maude's.
  \item Ms. Robbins's.
\end{choices}
\sectitle{第二节（共15小题：每小题1.5分，满分22.5分）}
\par 听下面 5 段对话或独白。每段对话或独白后有几个小题，从题中所给的 A、B、C 三个选项中选出最佳选项。听每段对话或独白前，你将有时间阅读各个小题，每小题 5 秒钟；听完后，各小题将给出 5 秒钟的作答时间。每段对话或独白读两遍。\par
\par 听下面一段对话，回答第6和第7题。\par
\question{6. What is the probable relationship between the speakers?}
\begin{choices}(3)
  \item Husband and wife.
  \item Uncle and nephew.
  \item Neighbors.
\end{choices}
\question{7. What will a cat do when it feels anxious?}
\begin{choices}(3)
  \item Hold its tail up.
  \item Stay close to the ground.
  \item Keep its eyes half-closed.
\end{choices}
\par 听第 7 段材料，回答第 8 和第 9 题。\par
\question{8. What did the man do before 2010?}
\begin{choices}(3)
  \item A manager.
  \item A teacher.
  \item A consultant.
\end{choices}
\question{9. What is the man's second suggestion?}
\begin{choices}(3)
  \item Setting a shared goal.
  \item Chatting with students.
  \item Holding a weekly class meeting.
\end{choices}
\par 听第 8 段材料，回答第 10 至 13 题。\par
\question{10. How does the woman sound at first?}
\begin{choices}(3)
  \item Grateful.
  \item Anxious.
  \item Doubtful.
\end{choices}
\question{11. What do we know about the word Hogmanay?}
\begin{choices}(3)
  \item It was invented in 1561.
  \item Its origin is still unclear.
  \item It was widely used in France.
\end{choices}
\question{12. What is the tradition of the first footing during Hogmanay?}
\begin{choices}(3)
  \item Bringing bread and coal as gifts.
  \item Visiting friends before midnight.
  \item Watching a traditional comedy with friends.
\end{choices}
\question{13. Where does the woman come from?}
\begin{choices}(3)
  \item Germany.
  \item France.
  \item Britain.
\end{choices}
\par 听第 9 段材料，回答第 14 至 16 题。\par
\question{14. What is the advantage of the shuttle bus compared to the Greyhound bus?}
\begin{choices}(3)
  \item Lower cost.
  \item Easier booking process.
  \item More frequent departures.
\end{choices}
\question{15. How much will the woman pay?}
\begin{choices}(3)
  \item \$15.
  \item \$28.
  \item \$35.
\end{choices}
\question{16. What will the woman do next?}
\begin{choices}(3)
  \item Catch a flight.
  \item Look for her credit card.
  \item Book a private hotel.
\end{choices}
\par 听第 10 段材料，回答第 17 至 20 题。\par
\question{17. What is the speaker doing?}
\begin{choices}(3)
  \item Reporting a piece of news.
  \item Providing first aid training.
  \item Hosting an award ceremony.
\end{choices}
\question{18. What symptom did Jessica notice in her client?}
\begin{choices}(3)
  \item The sound of choking.
  \item A sudden drop in blood pressure.
  \item Inability to communicate properly.
\end{choices}
\question{19. What was the immediate action taken by Jessica?}
\begin{choices}(2)
  \item She called 911 at once. 3. She updated the client's situation.
  \item She gave the client some medicine.
\end{choices}
\question{20. What is mainly introduced by the speaker?}
\begin{choices}(3)
  \item Jessica's lifesaving efforts.
  \item Jessica's chat with a client.
  \item Richard's race against time.
\end{choices}
\parttitle{第二部分 阅读理解（共两节，满分 50 分）}
\sectitle{第一节（共15小题：每小题2.5分，满分37.5分）}
\par 阅读下列短文，从每题所给的 A、B、C 和 D 四个选项中，选出最佳答案。\par
\passagelabel{A}
\par We've combed the globe for the best hotels of the year. Here are some of our picks. Revamped Hotel (Badeschloss, Austria)\par
\par Viennese architects BWM have transformed it into a red-hot retreat. Its 13-storey concrete addition has brought the room total to 102, and all are designed in candy colours with bathhouse accents like pool ladders, private saunas and deep tubs you can practically swim in. The icing on this cake is the rooftop heated pool atop three floors of saunas.\par
\par Remote Sleep (Boca de Agua, Mexico)\par
\par With a complex of 22 elevated treehouses deep in the rainforest on the Yucatán Peninsula, the treehouses have been created using wasted materials and a pared-back aesthetic that allows the surrounding jungle to take centre stage. The emphasis is on softening the walls between the wilderness and guests, with outdoor showers, views of hammocks swinging lazily on the deck.\par
\par Eco-Stay(Santo Mine, Greece)\par
\par In the whitewashed village of Oia, in northern Santorini, Santo Mine has been hot on sustainability. Supporting local suppliers, managing waste responsibly and protecting biodiversity with indigenous gardens is the MO (Mission Objective). Simple but elegant suites with light woods and sea-facing beds sleep up to four people with private pools, private whirlpool baths and the spectacular daily sunset.\par
\par Mountain Refuge (Odles Lodge, Italy)\par
\par Odles is a mix of understated luxury and raw wilderness, with pine forest and flower-flecked meadows at its feet. The minimalist-cool lodges use blonde spruce and Swiss stone pine — the fragrant untreated timber apparently promotes better sleep — while the windows and balconies frame impressive views. The spirit of South Tyrol's traditional mountain huts is captured, too.\par
\question{21. In Boca de Agua, travelers can \blanknum{}.}
\begin{choices}(2)
  \item soften the walls
  \item live in a treehouse
  \item perform on the centre stage
  \item view hammocks in the jungle
\end{choices}
\question{22. What is the main concern of the hotel in Greece?}
\begin{choices}(2)
  \item Eco-friendliness.
  \item Natural scenery.
  \item Customers' privacy.
  \item Social responsibility.
\end{choices}
\question{23. Which aspect is included in all of the four hotels mentioned above?}
\begin{choices}(2)
  \item Building materials.
  \item Specific location.
  \item Rooms available.
  \item Distinctive feature.
\end{choices}
\passagelabel{B}
\par Do we really need real dubbing artists in the age of artificial intelligence? Indeed, many have harbored doubts about the future and necessity of real dubbing(配音) artists. Yet for Kathryn Lin, a leading dubbing artist from China's Taiwan island, the profession means the shared happiness of generations of audiences on both sides of the Taiwan Strait.\par
\par Born in 1975, Lin was frequently chosen by her school to take part in various speech contests. Unfortunately, in her first year of college, her family faced a major financial crisis, and she had no choice but to drop out of college and take up part-time jobs in shops and restaurants. The chance to join the dubbing industry came when the dubbing artists' union on Taiwan island wanted to recruit students for dubbing performances, and she made it.\par
\par Since she was a newcomer to the industry, she had to complete her apprenticeship(学徒期) before getting real work, which usually takes three to five years of hard work, but Lin achieved that in less than two years.\par
\par "After working for a few years, I realized that the job was not just about bread and butter," she said, stressing that the popular characters she has dubbed for have given her "a lot of insight into life and intriguing experiences". Talking about the Chinese mainland's animation market, Lin said mainland productions have a lot of original elements. Mainland animations "are a cultural blend of the East and the West, and the style inspires me a lot when dubbing."\par
\par In her spare time, she also attends lectures to help young criminals, and children who have dropped out of school. "They might have been under stress or were really confused before breaking the law, or might have dropped out of school for family reasons," she said.\par
\par When it comes to AI dubbing, Lin said she is "actually trying to be less pessimistic...I think it's a future trend."\par
\par Lin believes dubbing artists can transform themselves by, for instance, teaching the art of dubbing to younger people, publishing books, teaching in schools and giving lectures.\par
\question{24. What made Lin take up dubbing?}
\begin{choices}(2)
  \item Her family's request.
  \item Her school's recommendation.
  \item The dubbing union's recruitment.
  \item Her keen interest in the dubbing industry.
\end{choices}
\question{25. What do we know about mainland productions in Lin's eyes?}
\begin{choices}(2)
  \item They are fully original.
  \item Their style can greatly inspire viewers.
  \item Their figures are appealing and thought-provoking.
  \item They feature a combination of the East and the West.
\end{choices}
\question{26. Which of the following best describes Lin?}
\begin{choices}(2)
  \item Diligent and thoughtful.
  \item Talented and generous.
  \item Competent and easy-going.
  \item Outgoing and humorous.
\end{choices}
\question{27. What did Lin think of dubbing in the future?}
\begin{choices}(2)
  \item She is pessimistic about AI dubbing.
  \item AI dubbing will substitute for dubbing artists to an extent.
  \item AI dubbing will merely make up a small part of the industry.
  \item Dubbing artists won't be replaced as dubbing represents shared happiness.
\end{choices}
\passagelabel{C}
\par You'll find it at the doctor's office, at the gym, and in online calculators: the body mass index, or BMI. It has been ever-present since its introduction in the 1970s. But this metric (度量标准) of a "healthy weight" has come under examination in the past decade or so, and with good reason. The relationship between physical well-being and BMI is complicated, and it's often not a good indicator of your overall health, experts say.\par
\par BMI is a simple measurement calculated by dividing a person's weight (or mass in kilograms if using metric) by their height. For decades, clinicians have used BMI to put our weight into different categories: underweight, normal, overweight, or obese. Twenty-five is the magic number: at this number or higher, you fall into the overweight or obese category.\par
\par Some population studies have associated having a high BMI with increased risk of cardiovascular disease, Type II diabetes, and other endocrine disorders. However, it may not be weight alone that plays into these negative health effects: research shows the societal shame people with high BMI have could also contribute to poor health outcomes.\par
\par Other studies have failed to find a strong correlation between BMI and health. People who are lean can be unhealthy, and people who have a BMI over 25 can have some health markers, says Fatima Stanford, an associate professor at Harvard Medical School and obesity medicine physician at Massachusetts General Hospital.\par
\par Then why do we still rely on BMI? One reason is that it's a relatively simple measurement to take. This makes the metric easily accessible to practitioners without specialized training, says Eleanna De Filippis, a physician at the Mayo Clinic. Stanford says that extremely low or high values of BMI could be an indicator of potential health issues despite its early introduction, but values around 25 to 30 are difficult to interpret without other health metrics.\par
\par However, there are big issues with only using BMI. to assess health. BMI can't tell what percentage of a person's weight is from their fat, muscle, or bone. That's why muscular athletes often have a high BMI despite having little body fat, De Filippis says. It's also possible to have poor health measures, like high blood pressure and cholesterol, and be at a "normal" BMI, Stanford says.\par
\question{28. What do we know about BMI?}
\begin{choices}(2)
  \item People within normal BMI range can be unhealthy.
  \item Experts suggest using BMI to imply our overall health.
  \item People who have a BMI of 26 fall into the obese category.
  \item People don't often use BMI to measure their health nowadays.
\end{choices}
\question{29. What does the underlined word "correlation" in paragraph 4 mean?}
\begin{choices}(2)
  \item Assumption.
  \item Connection.
  \item Correction.
  \item Distinction.
\end{choices}
\question{30. Why does Stanford think BMI is widely used?}
\begin{choices}(2)
  \item It is scientific and reliable in many cases.
  \item It is easy to use for those without much training.
  \item It is still an accurate indicator of your overall health.
  \item It could be of some use in indicating health problems.
\end{choices}
\question{31. What's the author's purpose in writing the text?}
\begin{choices}(2)
  \item To introduce to readers what BMI is.
  \item To argue that BMI should not be used widely.
  \item To inform readers of the shortcomings of BMI.
  \item To encourage readers to abandon the use of BMI.
\end{choices}
\passagelabel{D}
\par Why are some of us good at math, or writing, while others excel at art or basketball? What causes the differences? Are they largely determined by the genes(nature) we inherit(遗传) from our parents? How much does the social environment (nurture) have to do with it? If you were a scientist interested in these questions, wouldn't you love to study identical twins?\par
\par Over the last 10 years, a team of scientists led by psychologist Thomas J. Bouchard Jr. has studied about 65 pairs of identical twins who were raised apart. The scientists bring each pair of twins to the University of Minnesota for a week of intensive testing. They record the twins' height, weight, eye color, ear shape, and head length. Meanwhile, they give the twins IQ and personality tests.\par
\par Bouchard's team has been startled by the similarities between twins raised apart. The twins often have surprisingly similar gestures and postures, for instance. And they have led remarkably similar lives. Adult identical twins also tend to have similar medical histories, developing the eye disease glaucoma at the same time, for example.\par
\par "The coincidences are not surprising," some critics of the Minnesota study say. They argue that everyone's life has so many details that a number of coincidences are bound to exist. What's more, for every coincidence discovered by a pair of identical twins raised apart, the critics could point to a vast number of undiscovered differences.\par
\par According to Bouchard, they find that genes play a big part in shaping our personalities—helping to determine whether we respect tradition and like to follow rules or not. The team has reported that intelligence also seems to be influenced much more by genes than by environment. Despite being raised by different families, separated identical twins earned identical or nearly identical scores on adult intelligence tests.\par
\par Other scientists, however, disagree. "How can you study whether intelligence is inherited when there are so many different ways just to define intelligence?" The problem applies to other traits (特征), they say.\par
\par Bouchard himself makes the point that even the most closely matched twins are different from each other. So even though genes may have a strong influence over our lives, they're not the only influence. Our day-to-day experiences help to shape us too.\par
\question{32. Why are scientists so interested in studying identical twins raised apart in different environments?}
\begin{choices}(2)
  \item Because it can be studied with a week of intensive testing.
  \item Because this group of twins have the most striking similarities.
  \item Because it is a good way to measure the effects of nature and nurture.
  \item Because there has already been a lot of research on identical twins raised together.
\end{choices}
\question{33. What do critics of the Minnesota study say about the similarities between twins raised apart?}
\begin{choices}(2)
  \item They are merely coincidences.
  \item They are overstated to some extent.
  \item They are nothing compared with the differences.
  \item They are too common to deserve any serious consideration.
\end{choices}
\question{34. Which of the following statements would Bouchard agree to?}
\begin{choices}(2)
  \item Similarities go hand in hand with differences.
  \item Nature has a notable influence over our lives.
  \item Intelligence is something inherited rather than acquired.
  \item Environment plays a bigger part in personal development.
\end{choices}
\question{35. Which of the following is the best title for the text?}
\begin{choices}(2)
  \item Inherited Genes Favor Twins
  \item Environment Tells Apart Twins
  \item Science Looks Twice at Twins
  \item Physical Traits Come to Twins
\end{choices}
\sectitle{第二节（共5小题；每小题2.5分，满分12.5分）}
\par 阅读下面短文，从短文后的选项中选出能填入空白处的最佳选项。选项中有两项为多余选项。\par
\par Drinking coffee has been linked to a reduced risk of developing type II diabetes, neurological conditions such as dementia, and even death. \blanknum{36}\par
\par How decaf compares to regular coffee\par
\par The good news is that many of coffee's health benefits are still found in decaffeinated(脱咖啡因) coffee. The reason for this is thought to be due to the levels of antioxidants(抗氧化剂) in coffee beans, many of which remain even after the caffeine has been removed. Reduced though, the overall levels of antioxidants are still high. \blanknum{37}\par
\par The pros and cons of caffeine\par
\par Researchers found a reduced risk of cardiovascular disease and death in all coffee drinkers, which included people who drank decaf. Also, caffeine can help alleviate migraines (偏头痛) \blanknum{38} In the short-term, consuming too much caffeine can cause jitteriness, sleep disruptions, headaches, or a racing heart. Long-term, consuming too much caffeine can lead to withdrawal symptoms, such as headaches, if a person tries to stop.\par
\par One of the reasons that the health benefits of coffee have been debated for so long is due to the fact that its consumption can often go hand in hand with other habits. Researchers found coffee drinking was associated with a higher rate of smoking cigarettes — and that the health risks they faced such as lung cancer were due to smoking rather than coffee. Another study found, regular coffee consumption is associated with a reduced amount of weight gain over time. \blanknum{40} This held true whether people were drinking regular or decaf coffee. "Having unsweetened coffee seems to be the way to go," Ball says.\par
\begin{choices}(1)
  \item Context matters.
  \item The hidden cause and effect.
  \item The effects of caffeine can vary, however.
  \item That can translate to real benefits for decaf drinkers.
  \item Physicians should stop telling patients to stop drinking coffee.
  \item But only if drinkers had their coffee black or with cream, but no sugar.
  \item But as people increasingly turn to decaf, are they still getting those same benefits?
\end{choices}
\parttitle{第三部分 语言知识运用（共两节，满分30分）}
\sectitle{第一节（共15小题；每小题1分，满分15分）}
\par 阅读下面短文，从短文后各题所给的A、B、C和D四个选项中，选出可以填入空白处的最佳选项。\par
\par Jerry was the kind of guy who was always in a good mood and always had something positive to say. One day I went up to Jerry and asked him, "I don't \blanknum{41} it! You can't be in a good mood all of the time. How do you do it?" Jerry replied, "Each morning I wake up and say to myself, Jerry, you have two choices today — to be in a good mood or in a bad one. I choose the former." "Yeah, it's not that easy," I \blanknum{42}.\par
\par "Yes. it is," Jerry said. "Life is all about choices. You choose how you \blanknum{} 43\blanknum{} to situations. You choose how people will affect your mood. It's your choice how you live a life. "I \blanknum{} 44\blanknum{} what Jerry said. Soon thereafter, I changed my job. We lost touch. Several years later. I heard that Jerry was \blanknum{} 45\blanknum{} and was shot. I saw Jerry about six months after the accident.\par
\par I asked him what had \blanknum{46} his mind when he was taken to the hospital. Jerry replied, "I had two choices: I could choose to live, or I could choose to die. I chose to live." Jerry continued, "The nurses \blanknum{47} telling me I was going to be fine. But when they \blanknum{48} me into the emergency room, the expressions on the faces of the doctors and nurses \blanknum{49} me. In their eyes, I \blanknum{50}, ‘He’s a dead man.’ I knew I needed to take action." "What did you do?" I asked. "Well, before the \blanknum{51}, a nurse asked if I was allergic to(对……过敏) anything. The doctors and nurses were waiting for my reply... I took a deep breath and yelled. ‘\blanknum{52}!’ Over their laughter, I told them, ‘Operate on me as if I am alive, not dead’." Jerry would not have survived without the \blanknum{53} of his doctors, but also because of his amazing \blanknum{54}. I learned from him that every day we have the choice to live \blanknum{55} thanks to Jerry.\par
\cz{41}{make}{get}{keep}{beat}
\cz{42}{promised}{admitted}{protested}{commented}
\cz{43}{react}{attend}{reply}{devote}
\cz{44}{carried on}{depended on}{reflected on}{held on}
\cz{45}{swallowed}{sentenced}{accused}{robbed}
\cz{46}{gone through}{broken through}{pulled through}{carried through}
\cz{47}{kept}{risked}{regretted}{forgot}
\cz{48}{involved}{signaled}{wheeled}{persuaded}
\cz{49}{inspired}{reminded}{transformed}{scared}
\cz{50}{claimed}{read}{agreed}{pronounced}
\cz{51}{survival}{struggle}{surgery}{strategy}
\cz{52}{Hospital}{Mango}{Alcohol}{Bullets}
\cz{53}{ambition}{expertise}{sympathy}{enthusiasm}
\cz{54}{humor}{talent}{patience}{attitude}
\cz{55}{comfortably}{optimistically}{peacefully}{simply}
\sectitle{第二节（共10小题；每小题1.5分，满分15分）}
\par 阅读下面短文，在空白处填入1个适当的单词或括号内单词的正确形式。\par
\par The Winter Solstice, also known as the winter festival, is one of the 24 solar terms in China which falls on December 22, or one day before or after it. On that day, the sun is directly on the Tropic of Capricorn, which means the day in the northern hemisphere is the \blanknum{56} (short) day and the night longest. As the Winter Solstice is the first term \blanknum{57}(establish), it is often regarded as the starting point of China's 24 solar terms. Since this day \blanknum{58}(mark)a major change in climate, the ancient people attached great importance to it, \blanknum{59}(view) it as significant as the Spring Festival.\par
\par The Winter Solstice, a festival with \blanknum{60} history of more than 2,000 years originated from the Han Dynasty and reached its peak in glory in the Tang and Song Dynasties. In the Tang Dynasty, to celebrate the festival, the emperor and his ministers offered sacrifices \blanknum{61} heaven and held court meetings, while the common people held celebration \blanknum{62} (ceremony) to welcome the upcoming year.\par
\par Basically, the celebration of Winter Solstice nowadays \blanknum{63} (reflect) in different eating customs in different parts of China. In the north, people tend to eat dumplings, \blanknum{64} shape looks like ears, in order to prevent people from frostbite(冻疮) on ears, \blanknum{65} people in southern China prefer sweet dumplings, rice balls and long noodles to express their good wishes.\par
\parttitle{第四部分 写作（共两节，满分 40 分）}
\sectitle{第一节（满分15分）}
\par 假如你是李华，运动会临近，你想邀请外教 Mr. Smith 参加你们班的集体活动跳长绳(long rope jumping)请你写一封邮件给他，内容包括:(1)介绍该运动；(2)准备事宜。注意：\par
\par (1) 写作词数应为 80 左右;\par
\par (2) 请按如下格式在答题纸的相应位置作答。\par
\sectitle{第二节（满分25分）}
\par 阅读下面材料，根据其内容和所给段落开头语续写两段，使之构成一篇完整的短文。\par
\par I sat. He glared at me. "You can give it to me or I can take it, your choice. What will it be?" I pulled the package of hair ribbons out of the pocket of my jeans, handed it to him and begged, "You're not going to call my dad, are you?"\par
\par "I'm calling the police. They will call your father."\par
\par My head dropped onto my hands and I sobbed, "No, please! Can't you just let me go? I can pay you. I have money in my pocket. I'm only fourteen years old. Please, I won't ever shoplift again!"\par
\par "Save your tears. They won't work on me. I'm sick of you bad kids stealing, just for the thrill of it."\par
\par I sat there, trembling uncontrollably. The previous excitement brought by stealing lipstick, teen magazine and all that stuff... now gave way to immense fear and shame.\par
\par The police arrived, and the policemen walked me to their black-and-white car and opened the backdoor. I got in, and they drove me through the middle of our small town. I slouched down into the seat so no one could see me as I looked out the window at the evening sky. Then I saw the steeple(尖塔) of my family's church, and the guilt pierced me like a dagger. "How could I have been so stupid? I've broken my father's heart... and God's," thought I.\par
\par We arrived at the station, and a round woman with a square face asked me questions until I ran out of answers. She pointed to the door of a large open cell and said, "Sit. Wait." I walked in, and my footsteps made an echo that bounced off the bars. The tears started again as I sat down on a hard bench and heard her dial the telephone and say, "I have your daughter in a cell at the police station. No, she's not hurt. She was caught shopping. Can you come and get her? Okay. You're welcome, good-bye." She yelled, "Hey kid, your father's on his way."\par
\par 注意：(1) 续写词数应为 150 左右；(2) 请按如下格式在答题卡的相应位置作答。\par
\writingpara{Paragraph 1: About one hundred years later; I heard his voice say my name.}
\writinglines{8}
\writingpara{Paragraph 2: After we arrived home, Dad finally asked me, "Why?"}
\writinglines{8}
\end{document}
